\section{Exercise 3}

Mảng 10 phần tử và ba thanh ghi lưu trạng thái: \texttt{\$s1} là giá trị lớn nhất, \texttt{\$s2} là giá trị lớn thứ hai, \texttt{\$s3} là cờ báo đã có giá trị lớn thứ hai hay chưa. Phần tử đầu tiên được lưu vào mảng và cũng làm mốc khởi tạo cho \texttt{\$s1}:

\begin{codeblock}[language=MIPS]
  la $s0, arr               # $s0 = địa chỉ mảng
  li $s1, 0                 # $s1 = giá trị lớn nhất
  li $s2, 0                 # $s2 = giá trị lớn thứ hai
  li $s3, 0                 # $s3 = 1 khi đã có giá trị lớn thứ hai
  li $t0, 0                 # $t0 = chỉ số i

input_loop:
  beq $t0, 10, print_second # đủ 10 phần tử thì in kết quả

  li $v0, 4                 # syscall 4: in chuỗi
  la $a0, prompt1           # in "Please insert a element: "
  syscall

  li $v0, 5                 # syscall 5: đọc số nguyên
  syscall

  sll $t1, $t0, 2           # $t1 = 4i
  add $t1, $t1, $s0         # $t1 = base + 4i
  sw $v0, 0($t1)            # lưu vào phần tử thứ i

  beqz $t0, first_input     # phần tử đầu chưa có mốc so sánh
\end{codeblock}

Sau khi đã có phần tử đầu, mỗi số đọc vào rơi vào bốn trường hợp dưới đây. Số bằng \texttt{\$s1} bị bỏ qua vì nó cũng là giá trị lớn nhất, không phải giá trị lớn thứ hai:

\begin{codeblock}[language=MIPS]
  bgt $v0, $s1, new_max    # lớn hơn max: max cũ thành second
  beq $v0, $s1, next       # bằng max: bỏ qua
  beqz $s3, new_second     # chưa có second: nhận luôn
  bgt $v0, $s2, new_second # lớn hơn second: thay second
  j next

first_input:
  move $s1, $v0            # phần tử đầu làm max
  j next

new_max:
  move $s2, $s1            # second = max cũ
  move $s1, $v0            # max = phần tử mới
  li $s3, 1                # đã có second
  j next

new_second:
  move $s2, $v0            # second = phần tử mới
  li $s3, 1                # đã có second
  j next
\end{codeblock}

Chương trình in giá trị lớn thứ hai trước, sau đó là một vòng quét thứ hai tìm mọi chỉ số có phần tử bằng \texttt{\$s2}. Trong vòng quét đó, \texttt{\$t1} là cờ cho biết đã in chỉ số nào chưa: nó khởi tạo bằng 0, chỉ số khớp đầu tiên được in một mình rồi cờ bật lên 1; từ chỉ số thứ hai trở đi mới in thêm \texttt{", "} phía trước. Nhờ vậy kết quả không có dấu phẩy thừa ở đầu:

\begin{codeblock}[language=MIPS]
print_second:
  beqz $s3, no_second     # chưa có second: báo không tồn tại

  li $v0, 4               # syscall 4: in chuỗi
  la $a0, msg1            # chuỗi "Second largest value is "
  syscall

  li $v0, 1               # syscall 1: in số nguyên
  move $a0, $s2           # in giá trị lớn thứ hai
  syscall

  li $v0, 4               # syscall 4: in chuỗi
  la $a0, msg2            # chuỗi ", found in index "
  syscall

  li $t0, 0               # $t0 = chỉ số i
  li $t1, 0               # $t1 = 1 khi đã in ít nhất một chỉ số

find_second_indexes:
  beq $t0, 10, exit       # quét hết mảng thì kết thúc

  sll $t2, $t0, 2         # $t2 = 4i
  add $t2, $t2, $s0       # $t2 = base + 4i
  lw $t3, 0($t2)          # $t3 = phần tử thứ i

  bne $t3, $s2, find_next # không bằng second: bỏ qua
  beqz $t1, print_idx     # chỉ số đầu tiên: in không dấu phẩy

  li $v0, 4               # syscall 4: in chuỗi
  la $a0, msg3            # in ", " trước các chỉ số sau
  syscall

print_idx:
  li $v0, 1               # syscall 1: in số nguyên
  move $a0, $t0           # in chỉ số
  li $t1, 1               # đánh dấu đã in ít nhất một chỉ số
  syscall

find_next:
  addi $t0, $t0, 1        # i++
  j find_second_indexes   # tiếp tục quét
\end{codeblock}

Khi cờ \texttt{\$s3} vẫn là 0, nhãn \texttt{no\_second} in thông báo rồi thoát:

\begin{codeblock}[language=MIPS]
no_second:
  li $v0, 4    # syscall 4: in chuỗi
  la $a0, msg5 # "No second largest value exists."
  syscall
  j exit
\end{codeblock}

Nhập $1, 2, 7, -7, 3, 7, 4, 6, 8, 7$:

\begin{codeblock}[language={}]
Second largest value is 7, found in index 2, 5, 9.
\end{codeblock}

Với mảng toàn số bằng nhau, cờ \texttt{\$s3} không bao giờ được bật. Nhập $5, 5, 5, 5, 5, 5, 5, 5, 5, 5$:

\begin{codeblock}[language={}]
No second largest value exists.
\end{codeblock}
